%%%%%%%%%%%%%%%%%%%%%%%%%%%%%%%%%%%%%%%%%%%%%%%%%%%%%%%%%%%%
% 7.4 CONTINUOUS DISTRIBUTIONS
% The Composition of Advanced Mathematics
%%%%%%%%%%%%%%%%%%%%%%%%%%%%%%%%%%%%%%%%%%%%%%%%%%%%%%%%%%%%

\section{Continuous Distributions}
\label{sec:continuous-distributions}

\prerequisite{
Elementary probability, random variables, cumulative distribution
functions, derivatives, definite integrals, improper integrals,
exponential functions, logarithms, and basic algebra.
}

\learningobjectives{
By the end of this section, the reader should be able to:
\begin{itemize}
    \item define an absolutely continuous random variable;
    \item recognize and verify probability density functions;
    \item compute probabilities using definite integrals;
    \item derive cumulative distribution functions from densities;
    \item recover densities from differentiable CDFs;
    \item explain why a continuous random variable assigns probability
          zero to individual points;
    \item work with uniform distributions;
    \item work with exponential distributions;
    \item use the memoryless property of the exponential distribution;
    \item work with normal distributions;
    \item standardize normal random variables;
    \item use the standard normal CDF;
    \item understand location-scale normal families;
    \item recognize gamma distributions;
    \item recognize chi-square distributions as gamma special cases;
    \item recognize beta distributions;
    \item recognize Cauchy distributions;
    \item recognize Weibull distributions;
    \item understand Pareto distributions and heavy tails;
    \item work with triangular distributions;
    \item distinguish density, probability, and cumulative probability;
    \item compute quantiles for continuous distributions;
    \item select an appropriate continuous probability model;
    \item understand important relationships among standard continuous
          distributions.
\end{itemize}
}

%%%%%%%%%%%%%%%%%%%%%%%%%%%%%%%%%%%%%%%%%%%%%%%%%%%%%%%%%%%%
\subsection{What Is a Continuous Distribution?}
\label{subsec:what-is-continuous-distribution}
%%%%%%%%%%%%%%%%%%%%%%%%%%%%%%%%%%%%%%%%%%%%%%%%%%%%%%%%%%%%

A continuous probability distribution describes a random variable whose
probability is spread across intervals of real numbers rather than
concentrated at isolated points.

Typical examples include:

\[
\text{time},
\qquad
\text{length},
\qquad
\text{mass},
\qquad
\text{temperature},
\qquad
\text{voltage}.
\]

The central object is the probability density function.

%%%%%%%%%%%%%%%%%%%%%%%%%%%%%%%%%%%%%%%%%%%%%%%%%%%%%%%%%%%%
\subsection{Probability Density Functions}
\label{subsec:probability-density-functions}
%%%%%%%%%%%%%%%%%%%%%%%%%%%%%%%%%%%%%%%%%%%%%%%%%%%%%%%%%%%%

\begin{definitionbox}{Probability Density Function}
A function \(f_X\) is a probability density function for a random
variable \(X\) if

\[
\boxed{
f_X(x)\geq0
}
\]

for all \(x\), and

\[
\boxed{
\int_{-\infty}^{\infty}
f_X(x)\,dx
=
1.
}
\]

Probabilities are then computed by

\[
\boxed{
P(a<X<b)
=
\int_a^b f_X(x)\,dx.
}
\]
\end{definitionbox}

The abbreviation

\[
\boxed{
\text{PDF}
}
\]

stands for probability density function.

%%%%%%%%%%%%%%%%%%%%%%%%%%%%%%%%%%%%%%%%%%%%%%%%%%%%%%%%%%%%
\subsection{Density Is Not Probability}
\label{subsec:density-not-probability-continuous}
%%%%%%%%%%%%%%%%%%%%%%%%%%%%%%%%%%%%%%%%%%%%%%%%%%%%%%%%%%%%

The value

\[
f_X(x)
\]

is a density, not a probability.

It may even exceed \(1\).

Only areas under the density curve represent probabilities.

Thus,

\[
\boxed{
\text{probability}
=
\text{area under the density}.
}
\]

%%%%%%%%%%%%%%%%%%%%%%%%%%%%%%%%%%%%%%%%%%%%%%%%%%%%%%%%%%%%
\subsection{Point Probabilities}
\label{subsec:continuous-point-probabilities}
%%%%%%%%%%%%%%%%%%%%%%%%%%%%%%%%%%%%%%%%%%%%%%%%%%%%%%%%%%%%

For an absolutely continuous random variable,

\[
\boxed{
P(X=x)=0
}
\]

for every real number \(x\).

Indeed,

\[
P(X=x)
=
\int_x^x f_X(t)\,dt
=
0.
\]

Therefore individual points carry no probability mass.

%%%%%%%%%%%%%%%%%%%%%%%%%%%%%%%%%%%%%%%%%%%%%%%%%%%%%%%%%%%%
\subsection{Endpoint Inclusion}
\label{subsec:continuous-endpoint-inclusion}
%%%%%%%%%%%%%%%%%%%%%%%%%%%%%%%%%%%%%%%%%%%%%%%%%%%%%%%%%%%%

Because point probabilities vanish,

\[
P(a<X<b)
=
P(a\leq X<b)
\]

and

\[
P(a<X\leq b)
=
P(a\leq X\leq b).
\]

Therefore,

\[
\boxed{
P(a<X<b)
=
P(a\leq X\leq b).
}
\]

For continuous distributions, endpoint inclusion does not affect
probability.

%%%%%%%%%%%%%%%%%%%%%%%%%%%%%%%%%%%%%%%%%%%%%%%%%%%%%%%%%%%%
\subsection{Verifying a Density}
\label{subsec:verifying-density}
%%%%%%%%%%%%%%%%%%%%%%%%%%%%%%%%%%%%%%%%%%%%%%%%%%%%%%%%%%%%

To verify a proposed PDF:

\begin{enumerate}
    \item check
    \[
    f(x)\geq0;
    \]

    \item compute
    \[
    \int_{-\infty}^{\infty}f(x)\,dx;
    \]

    \item verify the total integral equals \(1\).
\end{enumerate}

Both conditions are required.

%%%%%%%%%%%%%%%%%%%%%%%%%%%%%%%%%%%%%%%%%%%%%%%%%%%%%%%%%%%%
\subsection{Worked Example: Finding a Normalizing Constant}
\label{subsec:worked-density-normalization}
%%%%%%%%%%%%%%%%%%%%%%%%%%%%%%%%%%%%%%%%%%%%%%%%%%%%%%%%%%%%

\begin{workedexamplebox}{Normalizing a Linear Density}

Suppose

\[
f(x)
=
\begin{cases}
cx,&0\leq x\leq2,\\
0,&\text{otherwise}.
\end{cases}
\]

Find \(c\).

Normalization requires

\[
\int_0^2 cx\,dx=1.
\]

Thus,

\[
c
\left[
\frac{x^2}{2}
\right]_0^2
=
1.
\]

Therefore,

\[
2c=1.
\]

\begin{answerbox}
\[
\boxed{
c=\frac12.
}
\]
\end{answerbox}

\end{workedexamplebox}

%%%%%%%%%%%%%%%%%%%%%%%%%%%%%%%%%%%%%%%%%%%%%%%%%%%%%%%%%%%%
\subsection{Probability from a Density}
\label{subsec:probability-from-density}
%%%%%%%%%%%%%%%%%%%%%%%%%%%%%%%%%%%%%%%%%%%%%%%%%%%%%%%%%%%%

Suppose \(X\) has density \(f_X\).

Then

\[
\boxed{
P(a\leq X\leq b)
=
\int_a^b f_X(x)\,dx.
}
\]

If the support of \(X\) does not intersect part of the interval, that
portion contributes zero.

%%%%%%%%%%%%%%%%%%%%%%%%%%%%%%%%%%%%%%%%%%%%%%%%%%%%%%%%%%%%
\subsection{Worked Example: Interval Probability}
\label{subsec:worked-density-interval}
%%%%%%%%%%%%%%%%%%%%%%%%%%%%%%%%%%%%%%%%%%%%%%%%%%%%%%%%%%%%

\begin{workedexamplebox}{Probability Under a Linear Density}

Continue with

\[
f(x)
=
\begin{cases}
\dfrac{x}{2},&0\leq x\leq2,\\
0,&\text{otherwise}.
\end{cases}
\]

Find

\[
P(1\leq X\leq2).
\]

Compute

\[
P(1\leq X\leq2)
=
\int_1^2\frac{x}{2}\,dx.
\]

Thus,

\[
=
\left[
\frac{x^2}{4}
\right]_1^2
=
1-\frac14.
\]

\begin{answerbox}
\[
\boxed{
P(1\leq X\leq2)
=
\frac34.
}
\]
\end{answerbox}

\end{workedexamplebox}

%%%%%%%%%%%%%%%%%%%%%%%%%%%%%%%%%%%%%%%%%%%%%%%%%%%%%%%%%%%%
\subsection{Continuous Cumulative Distribution Functions}
\label{subsec:continuous-cumulative-distribution-functions}
%%%%%%%%%%%%%%%%%%%%%%%%%%%%%%%%%%%%%%%%%%%%%%%%%%%%%%%%%%%%

For a continuous random variable,

\[
\boxed{
F_X(x)
=
P(X\leq x).
}
\]

If \(X\) has density \(f_X\), then

\[
\boxed{
F_X(x)
=
\int_{-\infty}^{x}
f_X(t)\,dt.
}
\]

Thus the CDF accumulates density from the left.

%%%%%%%%%%%%%%%%%%%%%%%%%%%%%%%%%%%%%%%%%%%%%%%%%%%%%%%%%%%%
\subsection{Recovering a Density from a CDF}
\label{subsec:recover-density-from-cdf}
%%%%%%%%%%%%%%%%%%%%%%%%%%%%%%%%%%%%%%%%%%%%%%%%%%%%%%%%%%%%

If \(F_X\) is differentiable, then

\[
\boxed{
f_X(x)
=
F_X'(x).
}
\]

Therefore:

\[
\boxed{
\text{density}
=
\text{derivative of accumulated probability}.
}
\]

This relationship is fundamental throughout continuous probability.

%%%%%%%%%%%%%%%%%%%%%%%%%%%%%%%%%%%%%%%%%%%%%%%%%%%%%%%%%%%%
\subsection{Worked Example: CDF from a Density}
\label{subsec:worked-cdf-from-density}
%%%%%%%%%%%%%%%%%%%%%%%%%%%%%%%%%%%%%%%%%%%%%%%%%%%%%%%%%%%%

\begin{workedexamplebox}{CDF of a Linear Density}

Let

\[
f_X(x)
=
\begin{cases}
\dfrac{x}{2},&0\leq x\leq2,\\
0,&\text{otherwise}.
\end{cases}
\]

For

\[
x<0,
\]

\[
F_X(x)=0.
\]

For

\[
0\leq x\leq2,
\]

\[
F_X(x)
=
\int_0^x\frac{t}{2}\,dt
=
\frac{x^2}{4}.
\]

For

\[
x\geq2,
\]

\[
F_X(x)=1.
\]

\begin{answerbox}
\[
\boxed{
F_X(x)
=
\begin{cases}
0,&x<0,\\[1mm]
\dfrac{x^2}{4},&0\leq x\leq2,\\[2mm]
1,&x\geq2.
\end{cases}
}
\]
\end{answerbox}

\end{workedexamplebox}

%%%%%%%%%%%%%%%%%%%%%%%%%%%%%%%%%%%%%%%%%%%%%%%%%%%%%%%%%%%%
\subsection{Survival Function}
\label{subsec:continuous-survival-function}
%%%%%%%%%%%%%%%%%%%%%%%%%%%%%%%%%%%%%%%%%%%%%%%%%%%%%%%%%%%%

The survival function is

\[
\boxed{
S_X(x)
=
P(X>x).
}
\]

For any continuous random variable,

\[
\boxed{
S_X(x)
=
1-F_X(x).
}
\]

Survival functions are especially important for lifetime and reliability
models.

%%%%%%%%%%%%%%%%%%%%%%%%%%%%%%%%%%%%%%%%%%%%%%%%%%%%%%%%%%%%
\subsection{Hazard Function Preview}
\label{subsec:hazard-function-preview}
%%%%%%%%%%%%%%%%%%%%%%%%%%%%%%%%%%%%%%%%%%%%%%%%%%%%%%%%%%%%

For a nonnegative continuous lifetime variable \(X\), define the hazard
function

\[
\boxed{
h_X(x)
=
\frac{f_X(x)}{S_X(x)}
}
\]

whenever

\[
S_X(x)>0.
\]

Informally, \(h_X(x)\) measures the instantaneous rate of failure near
time \(x\), conditional on survival up to \(x\).

This concept is developed further in survival analysis.

%%%%%%%%%%%%%%%%%%%%%%%%%%%%%%%%%%%%%%%%%%%%%%%%%%%%%%%%%%%%
\subsection{Uniform Distribution}
\label{subsec:continuous-uniform-distribution}
%%%%%%%%%%%%%%%%%%%%%%%%%%%%%%%%%%%%%%%%%%%%%%%%%%%%%%%%%%%%

\begin{definitionbox}{Continuous Uniform Distribution}
A random variable \(X\) is uniformly distributed on

\[
[a,b]
\]

if

\[
\boxed{
f_X(x)
=
\begin{cases}
\dfrac{1}{b-a},&a\leq x\leq b,\\[2mm]
0,&\text{otherwise}.
\end{cases}
}
\]

We write

\[
\boxed{
X\sim\operatorname{Uniform}(a,b).
}
\]
\end{definitionbox}

Every interval of equal length inside \([a,b]\) has equal probability.

%%%%%%%%%%%%%%%%%%%%%%%%%%%%%%%%%%%%%%%%%%%%%%%%%%%%%%%%%%%%
\subsection{Uniform CDF}
\label{subsec:uniform-cdf}
%%%%%%%%%%%%%%%%%%%%%%%%%%%%%%%%%%%%%%%%%%%%%%%%%%%%%%%%%%%%

If

\[
X\sim\operatorname{Uniform}(a,b),
\]

then

\[
\boxed{
F_X(x)
=
\begin{cases}
0,&x<a,\\[1mm]
\dfrac{x-a}{b-a},&a\leq x\leq b,\\[3mm]
1,&x>b.
\end{cases}
}
\]

The CDF is linear across the support.

%%%%%%%%%%%%%%%%%%%%%%%%%%%%%%%%%%%%%%%%%%%%%%%%%%%%%%%%%%%%
\subsection{Uniform Interval Probabilities}
\label{subsec:uniform-interval-probabilities}
%%%%%%%%%%%%%%%%%%%%%%%%%%%%%%%%%%%%%%%%%%%%%%%%%%%%%%%%%%%%

For

\[
a\leq c<d\leq b,
\]

\[
P(c\leq X\leq d)
=
\int_c^d
\frac1{b-a}\,dx.
\]

Thus,

\[
\boxed{
P(c\leq X\leq d)
=
\frac{d-c}{b-a}.
}
\]

Probability equals the fraction of total interval length.

%%%%%%%%%%%%%%%%%%%%%%%%%%%%%%%%%%%%%%%%%%%%%%%%%%%%%%%%%%%%
\subsection{Worked Example: Uniform Waiting Time}
\label{subsec:worked-uniform-waiting}
%%%%%%%%%%%%%%%%%%%%%%%%%%%%%%%%%%%%%%%%%%%%%%%%%%%%%%%%%%%%

\begin{workedexamplebox}{Uniform Arrival Time}

Suppose an arrival time is uniformly distributed between \(0\) and
\(20\) minutes.

Let

\[
X\sim\operatorname{Uniform}(0,20).
\]

Find the probability the arrival occurs between \(5\) and \(12\)
minutes.

\[
P(5\leq X\leq12)
=
\frac{12-5}{20}.
\]

\begin{answerbox}
\[
\boxed{
P(5\leq X\leq12)
=
\frac7{20}.
}
\]
\end{answerbox}

\end{workedexamplebox}

%%%%%%%%%%%%%%%%%%%%%%%%%%%%%%%%%%%%%%%%%%%%%%%%%%%%%%%%%%%%
\subsection{Exponential Distribution}
\label{subsec:exponential-distribution}
%%%%%%%%%%%%%%%%%%%%%%%%%%%%%%%%%%%%%%%%%%%%%%%%%%%%%%%%%%%%

The exponential distribution is a fundamental model for positive
waiting times.

\begin{definitionbox}{Exponential Distribution}
A random variable \(X\) has an exponential distribution with rate
\(\lambda>0\), written

\[
\boxed{
X\sim\operatorname{Exp}(\lambda),
}
\]

if

\[
\boxed{
f_X(x)
=
\begin{cases}
\lambda e^{-\lambda x},&x\geq0,\\
0,&x<0.
\end{cases}
}
\]
\end{definitionbox}

The parameter

\[
\lambda
\]

is called the rate.

%%%%%%%%%%%%%%%%%%%%%%%%%%%%%%%%%%%%%%%%%%%%%%%%%%%%%%%%%%%%
\subsection{Exponential Normalization}
\label{subsec:exponential-normalization}
%%%%%%%%%%%%%%%%%%%%%%%%%%%%%%%%%%%%%%%%%%%%%%%%%%%%%%%%%%%%

Check:

\[
\int_0^\infty
\lambda e^{-\lambda x}\,dx.
\]

We have

\[
=
\left[
-e^{-\lambda x}
\right]_0^\infty.
\]

Therefore,

\[
\boxed{
\int_0^\infty
\lambda e^{-\lambda x}\,dx
=
1.
}
\]

So the exponential density is properly normalized.

%%%%%%%%%%%%%%%%%%%%%%%%%%%%%%%%%%%%%%%%%%%%%%%%%%%%%%%%%%%%
\subsection{Exponential CDF}
\label{subsec:exponential-cdf}
%%%%%%%%%%%%%%%%%%%%%%%%%%%%%%%%%%%%%%%%%%%%%%%%%%%%%%%%%%%%

For

\[
x\geq0,
\]

\[
F_X(x)
=
\int_0^x
\lambda e^{-\lambda t}\,dt.
\]

Thus,

\[
\boxed{
F_X(x)
=
1-e^{-\lambda x}.
}
\]

The complete CDF is

\[
\boxed{
F_X(x)
=
\begin{cases}
0,&x<0,\\
1-e^{-\lambda x},&x\geq0.
\end{cases}
}
\]

%%%%%%%%%%%%%%%%%%%%%%%%%%%%%%%%%%%%%%%%%%%%%%%%%%%%%%%%%%%%
\subsection{Exponential Survival Function}
\label{subsec:exponential-survival}
%%%%%%%%%%%%%%%%%%%%%%%%%%%%%%%%%%%%%%%%%%%%%%%%%%%%%%%%%%%%

Since

\[
S_X(x)=1-F_X(x),
\]

we obtain

\[
\boxed{
S_X(x)
=
e^{-\lambda x},
\qquad
x\geq0.
}
\]

Therefore,

\[
\boxed{
P(X>x)
=
e^{-\lambda x}.
}
\]

%%%%%%%%%%%%%%%%%%%%%%%%%%%%%%%%%%%%%%%%%%%%%%%%%%%%%%%%%%%%
\subsection{Worked Example: Exponential Waiting Time}
\label{subsec:worked-exponential-wait}
%%%%%%%%%%%%%%%%%%%%%%%%%%%%%%%%%%%%%%%%%%%%%%%%%%%%%%%%%%%%

\begin{workedexamplebox}{Waiting More Than Ten Minutes}

Suppose

\[
X\sim\operatorname{Exp}(0.2).
\]

Find

\[
P(X>10).
\]

Using the survival function,

\[
P(X>10)
=
e^{-(0.2)(10)}.
\]

\begin{answerbox}
\[
\boxed{
P(X>10)=e^{-2}.
}
\]
\end{answerbox}

\end{workedexamplebox}

%%%%%%%%%%%%%%%%%%%%%%%%%%%%%%%%%%%%%%%%%%%%%%%%%%%%%%%%%%%%
\subsection{Memoryless Property of the Exponential Distribution}
\label{subsec:exponential-memoryless}
%%%%%%%%%%%%%%%%%%%%%%%%%%%%%%%%%%%%%%%%%%%%%%%%%%%%%%%%%%%%

\begin{theorem}[Exponential Memoryless Property]
If

\[
X\sim\operatorname{Exp}(\lambda),
\]

then for

\[
s,t\geq0,
\]

\[
\boxed{
P(X>s+t\mid X>s)
=
P(X>t).
}
\]
\end{theorem}

Thus elapsed waiting time does not affect the remaining waiting-time
distribution.

%%%%%%%%%%%%%%%%%%%%%%%%%%%%%%%%%%%%%%%%%%%%%%%%%%%%%%%%%%%%
\subsection{Proof of Exponential Memorylessness}
\label{subsec:proof-exponential-memoryless}
%%%%%%%%%%%%%%%%%%%%%%%%%%%%%%%%%%%%%%%%%%%%%%%%%%%%%%%%%%%%

Because

\[
\{X>s+t\}
\subseteq
\{X>s\},
\]

we have

\[
P(X>s+t\mid X>s)
=
\frac{P(X>s+t)}{P(X>s)}.
\]

Using

\[
P(X>x)=e^{-\lambda x},
\]

we obtain

\[
=
\frac{
e^{-\lambda(s+t)}
}{
e^{-\lambda s}
}
=
e^{-\lambda t}.
\]

Therefore,

\[
\boxed{
P(X>s+t\mid X>s)
=
P(X>t).
}
\]

%%%%%%%%%%%%%%%%%%%%%%%%%%%%%%%%%%%%%%%%%%%%%%%%%%%%%%%%%%%%
\subsection{Exponential Hazard Rate}
\label{subsec:exponential-hazard}
%%%%%%%%%%%%%%%%%%%%%%%%%%%%%%%%%%%%%%%%%%%%%%%%%%%%%%%%%%%%

For the exponential distribution,

\[
f_X(x)
=
\lambda e^{-\lambda x}
\]

and

\[
S_X(x)
=
e^{-\lambda x}.
\]

Therefore,

\[
h_X(x)
=
\frac{
\lambda e^{-\lambda x}
}{
e^{-\lambda x}
}.
\]

Thus,

\[
\boxed{
h_X(x)=\lambda.
}
\]

The exponential distribution has constant hazard rate.

%%%%%%%%%%%%%%%%%%%%%%%%%%%%%%%%%%%%%%%%%%%%%%%%%%%%%%%%%%%%
\subsection{Exponential and Poisson Connection}
\label{subsec:exponential-poisson-connection}
%%%%%%%%%%%%%%%%%%%%%%%%%%%%%%%%%%%%%%%%%%%%%%%%%%%%%%%%%%%%

In a Poisson process with event rate \(\lambda\):

\[
\boxed{
\text{number of events in time }t
\sim
\operatorname{Poisson}(\lambda t),
}
\]

while the waiting time until the next event satisfies

\[
\boxed{
X\sim\operatorname{Exp}(\lambda).
}
\]

Thus Poisson counts and exponential waiting times describe two aspects of
the same stochastic model.

%%%%%%%%%%%%%%%%%%%%%%%%%%%%%%%%%%%%%%%%%%%%%%%%%%%%%%%%%%%%
\subsection{Normal Distribution}
\label{subsec:normal-distribution}
%%%%%%%%%%%%%%%%%%%%%%%%%%%%%%%%%%%%%%%%%%%%%%%%%%%%%%%%%%%%

The normal distribution is one of the most important distributions in
probability and statistics.

\begin{definitionbox}{Normal Distribution}
A random variable \(X\) is normally distributed with parameters
\(\mu\in\mathbb{R}\) and \(\sigma^2>0\), written

\[
\boxed{
X\sim N(\mu,\sigma^2),
}
\]

if

\[
\boxed{
f_X(x)
=
\frac{1}{\sigma\sqrt{2\pi}}
\exp
\left(
-\frac{(x-\mu)^2}{2\sigma^2}
\right),
\qquad
-\infty<x<\infty.
}
\]
\end{definitionbox}

The Greek letter

\[
\mu
\]

is called ``mu.''

The Greek letter

\[
\sigma
\]

is called ``sigma.''

%%%%%%%%%%%%%%%%%%%%%%%%%%%%%%%%%%%%%%%%%%%%%%%%%%%%%%%%%%%%
\subsection{Shape of the Normal Distribution}
\label{subsec:normal-shape}
%%%%%%%%%%%%%%%%%%%%%%%%%%%%%%%%%%%%%%%%%%%%%%%%%%%%%%%%%%%%

The normal density is:

\begin{itemize}
    \item bell-shaped;
    \item symmetric about \(\mu\);
    \item positive for every real \(x\);
    \item decreasing into both tails;
    \item completely determined by \(\mu\) and \(\sigma\).
\end{itemize}

The parameter \(\mu\) controls location.

The parameter \(\sigma\) controls spread.

%%%%%%%%%%%%%%%%%%%%%%%%%%%%%%%%%%%%%%%%%%%%%%%%%%%%%%%%%%%%
\subsection{Location and Scale in the Normal Family}
\label{subsec:normal-location-scale}
%%%%%%%%%%%%%%%%%%%%%%%%%%%%%%%%%%%%%%%%%%%%%%%%%%%%%%%%%%%%

Changing \(\mu\) shifts the distribution horizontally.

Increasing \(\sigma\) spreads the density over a wider range.

Decreasing \(\sigma\) concentrates probability more tightly around
\(\mu\).

Thus,

\[
\boxed{
\mu
=
\text{center},
\qquad
\sigma
=
\text{scale}.
}
\]

%%%%%%%%%%%%%%%%%%%%%%%%%%%%%%%%%%%%%%%%%%%%%%%%%%%%%%%%%%%%
\subsection{Standard Normal Distribution}
\label{subsec:standard-normal-distribution}
%%%%%%%%%%%%%%%%%%%%%%%%%%%%%%%%%%%%%%%%%%%%%%%%%%%%%%%%%%%%

The standard normal distribution has parameters

\[
\mu=0
\]

and

\[
\sigma^2=1.
\]

It is written

\[
\boxed{
Z\sim N(0,1).
}
\]

Its density is

\[
\boxed{
\phi(z)
=
\frac{1}{\sqrt{2\pi}}
e^{-z^2/2}.
}
\]

The Greek letter

\[
\phi
\]

is called ``phi.''

%%%%%%%%%%%%%%%%%%%%%%%%%%%%%%%%%%%%%%%%%%%%%%%%%%%%%%%%%%%%
\subsection{Standard Normal CDF}
\label{subsec:standard-normal-cdf}
%%%%%%%%%%%%%%%%%%%%%%%%%%%%%%%%%%%%%%%%%%%%%%%%%%%%%%%%%%%%

The standard normal CDF is conventionally denoted

\[
\boxed{
\Phi(z)
=
P(Z\leq z).
}
\]

The capital Greek letter

\[
\Phi
\]

is called ``capital phi.''

Thus,

\[
\boxed{
\Phi(z)
=
\int_{-\infty}^{z}
\frac{1}{\sqrt{2\pi}}
e^{-t^2/2}\,dt.
}
\]

This integral has no elementary antiderivative.

%%%%%%%%%%%%%%%%%%%%%%%%%%%%%%%%%%%%%%%%%%%%%%%%%%%%%%%%%%%%
\subsection{Normal Standardization}
\label{subsec:normal-standardization}
%%%%%%%%%%%%%%%%%%%%%%%%%%%%%%%%%%%%%%%%%%%%%%%%%%%%%%%%%%%%

If

\[
X\sim N(\mu,\sigma^2),
\]

define

\[
\boxed{
Z
=
\frac{X-\mu}{\sigma}.
}
\]

Then

\[
\boxed{
Z\sim N(0,1).
}
\]

This process is called standardization.

The quantity

\[
\frac{x-\mu}{\sigma}
\]

is commonly called a \(z\)-score.

%%%%%%%%%%%%%%%%%%%%%%%%%%%%%%%%%%%%%%%%%%%%%%%%%%%%%%%%%%%%
\subsection{Interpreting a \(z\)-Score}
\label{subsec:z-score-interpretation}
%%%%%%%%%%%%%%%%%%%%%%%%%%%%%%%%%%%%%%%%%%%%%%%%%%%%%%%%%%%%

A \(z\)-score measures how many standard deviations a value lies from
the center.

If

\[
z=2,
\]

the value is two standard deviations above \(\mu\).

If

\[
z=-1.5,
\]

the value is \(1.5\) standard deviations below \(\mu\).

%%%%%%%%%%%%%%%%%%%%%%%%%%%%%%%%%%%%%%%%%%%%%%%%%%%%%%%%%%%%
\subsection{Normal Probability Formula}
\label{subsec:normal-probability-formula}
%%%%%%%%%%%%%%%%%%%%%%%%%%%%%%%%%%%%%%%%%%%%%%%%%%%%%%%%%%%%

If

\[
X\sim N(\mu,\sigma^2),
\]

then

\[
P(X\leq x)
=
P
\left(
Z\leq
\frac{x-\mu}{\sigma}
\right).
\]

Therefore,

\[
\boxed{
P(X\leq x)
=
\Phi
\left(
\frac{x-\mu}{\sigma}
\right).
}
\]

%%%%%%%%%%%%%%%%%%%%%%%%%%%%%%%%%%%%%%%%%%%%%%%%%%%%%%%%%%%%
\subsection{Worked Example: Standardizing a Normal Variable}
\label{subsec:worked-normal-standardization}
%%%%%%%%%%%%%%%%%%%%%%%%%%%%%%%%%%%%%%%%%%%%%%%%%%%%%%%%%%%%

\begin{workedexamplebox}{Finding a Normal Probability}

Suppose

\[
X\sim N(100,15^2).
\]

Find

\[
P(X\leq130).
\]

Standardize:

\[
Z
=
\frac{130-100}{15}
=
2.
\]

Therefore,

\begin{answerbox}
\[
\boxed{
P(X\leq130)
=
\Phi(2).
}
\]
\end{answerbox}

\end{workedexamplebox}

%%%%%%%%%%%%%%%%%%%%%%%%%%%%%%%%%%%%%%%%%%%%%%%%%%%%%%%%%%%%
\subsection{Normal Symmetry}
\label{subsec:normal-symmetry}
%%%%%%%%%%%%%%%%%%%%%%%%%%%%%%%%%%%%%%%%%%%%%%%%%%%%%%%%%%%%

The standard normal density satisfies

\[
\phi(-z)=\phi(z).
\]

Therefore,

\[
\boxed{
\Phi(-z)
=
1-\Phi(z).
}
\]

Equivalently,

\[
\boxed{
P(Z\geq z)
=
P(Z\leq-z).
}
\]

This symmetry simplifies many probability calculations.

%%%%%%%%%%%%%%%%%%%%%%%%%%%%%%%%%%%%%%%%%%%%%%%%%%%%%%%%%%%%
\subsection{Worked Example: Central Normal Probability}
\label{subsec:worked-central-normal}
%%%%%%%%%%%%%%%%%%%%%%%%%%%%%%%%%%%%%%%%%%%%%%%%%%%%%%%%%%%%

\begin{workedexamplebox}{Probability Within Two Standard Deviations}

For

\[
Z\sim N(0,1),
\]

consider

\[
P(-2\leq Z\leq2).
\]

Using the CDF,

\[
P(-2\leq Z\leq2)
=
\Phi(2)-\Phi(-2).
\]

Using symmetry,

\[
\Phi(-2)=1-\Phi(2).
\]

Therefore,

\begin{answerbox}
\[
\boxed{
P(-2\leq Z\leq2)
=
2\Phi(2)-1.
}
\]
\end{answerbox}

\end{workedexamplebox}

%%%%%%%%%%%%%%%%%%%%%%%%%%%%%%%%%%%%%%%%%%%%%%%%%%%%%%%%%%%%
\subsection{The 68--95--99.7 Rule}
\label{subsec:empirical-rule-normal}
%%%%%%%%%%%%%%%%%%%%%%%%%%%%%%%%%%%%%%%%%%%%%%%%%%%%%%%%%%%%

For a normal distribution, approximately:

\[
\boxed{
68\%
}
\]

of probability lies within one standard deviation of the mean,

\[
\boxed{
95\%
}
\]

lies within two standard deviations,

and

\[
\boxed{
99.7\%
}
\]

lies within three standard deviations.

This is often called the empirical rule.

%%%%%%%%%%%%%%%%%%%%%%%%%%%%%%%%%%%%%%%%%%%%%%%%%%%%%%%%%%%%
\subsection{Normal Quantiles}
\label{subsec:normal-quantiles}
%%%%%%%%%%%%%%%%%%%%%%%%%%%%%%%%%%%%%%%%%%%%%%%%%%%%%%%%%%%%

Let

\[
z_p
\]

satisfy

\[
\boxed{
\Phi(z_p)=p.
}
\]

Then \(z_p\) is the \(p\)-quantile of the standard normal distribution.

For

\[
X\sim N(\mu,\sigma^2),
\]

the corresponding quantile is

\[
\boxed{
x_p
=
\mu+\sigma z_p.
}
\]

%%%%%%%%%%%%%%%%%%%%%%%%%%%%%%%%%%%%%%%%%%%%%%%%%%%%%%%%%%%%
\subsection{Worked Example: Normal Percentile}
\label{subsec:worked-normal-percentile}
%%%%%%%%%%%%%%%%%%%%%%%%%%%%%%%%%%%%%%%%%%%%%%%%%%%%%%%%%%%%

\begin{workedexamplebox}{Converting a Standard-Normal Quantile}

Suppose

\[
X\sim N(50,10^2)
\]

and the desired standard-normal quantile is

\[
z_p=1.645.
\]

Then

\[
x_p
=
50+10(1.645).
\]

\begin{answerbox}
\[
\boxed{
x_p=66.45.
}
\]
\end{answerbox}

\end{workedexamplebox}

%%%%%%%%%%%%%%%%%%%%%%%%%%%%%%%%%%%%%%%%%%%%%%%%%%%%%%%%%%%%
\subsection{Normal Approximation Preview}
\label{subsec:normal-approximation-preview}
%%%%%%%%%%%%%%%%%%%%%%%%%%%%%%%%%%%%%%%%%%%%%%%%%%%%%%%%%%%%

Many distributions become approximately normal under suitable
conditions.

For example, a binomial random variable with sufficiently large

\[
np
\]

and

\[
n(1-p)
\]

can often be approximated using a normal distribution.

This phenomenon is explained systematically by central limit theorems
later in the chapter.

%%%%%%%%%%%%%%%%%%%%%%%%%%%%%%%%%%%%%%%%%%%%%%%%%%%%%%%%%%%%
\subsection{Continuity Correction Preview}
\label{subsec:continuity-correction-preview}
%%%%%%%%%%%%%%%%%%%%%%%%%%%%%%%%%%%%%%%%%%%%%%%%%%%%%%%%%%%%

When approximating an integer-valued distribution by a continuous
normal distribution, probability assigned to a single integer can be
represented by an interval of width \(1\).

For example,

\[
P(X\leq k)
\]

is often approximated by

\[
\boxed{
P(Y\leq k+0.5)
}
\]

for a continuous normal approximation \(Y\).

This adjustment is called a continuity correction.

%%%%%%%%%%%%%%%%%%%%%%%%%%%%%%%%%%%%%%%%%%%%%%%%%%%%%%%%%%%%
\subsection{Gamma Distribution}
\label{subsec:gamma-distribution}
%%%%%%%%%%%%%%%%%%%%%%%%%%%%%%%%%%%%%%%%%%%%%%%%%%%%%%%%%%%%

The gamma distribution generalizes the exponential distribution.

\begin{definitionbox}{Gamma Distribution}
A random variable \(X\) has a gamma distribution with shape
\(\alpha>0\) and rate \(\lambda>0\) if

\[
\boxed{
f_X(x)
=
\frac{
\lambda^\alpha
}{
\Gamma(\alpha)
}
x^{\alpha-1}
e^{-\lambda x},
\qquad
x>0.
}
\]
\end{definitionbox}

The Greek letter

\[
\alpha
\]

is called ``alpha.''

The function

\[
\Gamma
\]

is the gamma function.

%%%%%%%%%%%%%%%%%%%%%%%%%%%%%%%%%%%%%%%%%%%%%%%%%%%%%%%%%%%%
\subsection{Gamma Function}
\label{subsec:gamma-function-probability}
%%%%%%%%%%%%%%%%%%%%%%%%%%%%%%%%%%%%%%%%%%%%%%%%%%%%%%%%%%%%

The gamma function is defined by

\[
\boxed{
\Gamma(\alpha)
=
\int_0^\infty
x^{\alpha-1}e^{-x}\,dx,
\qquad
\alpha>0.
}
\]

It satisfies

\[
\boxed{
\Gamma(\alpha+1)
=
\alpha\Gamma(\alpha).
}
\]

For positive integers \(n\),

\[
\boxed{
\Gamma(n)
=
(n-1)!.
}
\]

Thus the gamma function extends the factorial.

%%%%%%%%%%%%%%%%%%%%%%%%%%%%%%%%%%%%%%%%%%%%%%%%%%%%%%%%%%%%
\subsection{Exponential as a Gamma Special Case}
\label{subsec:exponential-gamma-special-case}
%%%%%%%%%%%%%%%%%%%%%%%%%%%%%%%%%%%%%%%%%%%%%%%%%%%%%%%%%%%%

Set

\[
\alpha=1.
\]

Since

\[
\Gamma(1)=1,
\]

the gamma density becomes

\[
\lambda e^{-\lambda x}.
\]

Therefore,

\[
\boxed{
\operatorname{Exp}(\lambda)
=
\operatorname{Gamma}(1,\lambda)
}
\]

under the shape-rate parameterization.

%%%%%%%%%%%%%%%%%%%%%%%%%%%%%%%%%%%%%%%%%%%%%%%%%%%%%%%%%%%%
\subsection{Erlang Distribution}
\label{subsec:erlang-distribution}
%%%%%%%%%%%%%%%%%%%%%%%%%%%%%%%%%%%%%%%%%%%%%%%%%%%%%%%%%%%%

When the gamma shape parameter is a positive integer,

\[
\alpha=n,
\]

the distribution is often called an Erlang distribution.

It naturally models waiting time until the \(n\)-th event in a Poisson
process.

Thus:

\[
\boxed{
\text{exponential}
=
\text{waiting for first event},
}
\]

while

\[
\boxed{
\text{Erlang}
=
\text{waiting for the \(n\)-th event}.
}
\]

%%%%%%%%%%%%%%%%%%%%%%%%%%%%%%%%%%%%%%%%%%%%%%%%%%%%%%%%%%%%
\subsection{Chi-Square Distribution}
\label{subsec:chi-square-distribution}
%%%%%%%%%%%%%%%%%%%%%%%%%%%%%%%%%%%%%%%%%%%%%%%%%%%%%%%%%%%%

The chi-square distribution is an important gamma special case.

\begin{definitionbox}{Chi-Square Distribution}
A chi-square random variable with \(\nu>0\) degrees of freedom has
density

\[
\boxed{
f_X(x)
=
\frac{
1
}{
2^{\nu/2}
\Gamma(\nu/2)
}
x^{\nu/2-1}
e^{-x/2},
\qquad
x>0.
}
\]
\end{definitionbox}

The Greek letter

\[
\nu
\]

is called ``nu.''

We write

\[
\boxed{
X\sim\chi^2_\nu.
}
\]

%%%%%%%%%%%%%%%%%%%%%%%%%%%%%%%%%%%%%%%%%%%%%%%%%%%%%%%%%%%%
\subsection{Chi-Square as Gamma}
\label{subsec:chi-square-gamma}
%%%%%%%%%%%%%%%%%%%%%%%%%%%%%%%%%%%%%%%%%%%%%%%%%%%%%%%%%%%%

Under the shape-rate convention,

\[
\boxed{
\chi^2_\nu
=
\operatorname{Gamma}
\left(
\frac{\nu}{2},
\frac12
\right).
}
\]

Thus chi-square theory is contained inside the broader gamma family.

Chi-square distributions become central in mathematical statistics and
hypothesis testing.

%%%%%%%%%%%%%%%%%%%%%%%%%%%%%%%%%%%%%%%%%%%%%%%%%%%%%%%%%%%%
\subsection{Chi-Square from Standard Normals}
\label{subsec:chi-square-standard-normal}
%%%%%%%%%%%%%%%%%%%%%%%%%%%%%%%%%%%%%%%%%%%%%%%%%%%%%%%%%%%%

If

\[
Z_1,\ldots,Z_\nu
\]

are independent standard normal random variables, then

\[
\boxed{
Z_1^2+\cdots+Z_\nu^2
\sim
\chi^2_\nu.
}
\]

This important relationship connects normal distributions with
chi-square distributions.

%%%%%%%%%%%%%%%%%%%%%%%%%%%%%%%%%%%%%%%%%%%%%%%%%%%%%%%%%%%%
\subsection{Beta Distribution}
\label{subsec:beta-distribution}
%%%%%%%%%%%%%%%%%%%%%%%%%%%%%%%%%%%%%%%%%%%%%%%%%%%%%%%%%%%%

The beta distribution is supported on

\[
[0,1].
\]

It is especially useful for modeling random proportions and
probabilities.

\begin{definitionbox}{Beta Distribution}
A random variable \(X\) has a beta distribution with parameters
\(\alpha,\beta>0\) if

\[
\boxed{
f_X(x)
=
\frac{
x^{\alpha-1}(1-x)^{\beta-1}
}{
B(\alpha,\beta)
},
\qquad
0<x<1.
}
\]
\end{definitionbox}

The Greek letter

\[
\beta
\]

is called ``beta.''

%%%%%%%%%%%%%%%%%%%%%%%%%%%%%%%%%%%%%%%%%%%%%%%%%%%%%%%%%%%%
\subsection{Beta Function}
\label{subsec:beta-function}
%%%%%%%%%%%%%%%%%%%%%%%%%%%%%%%%%%%%%%%%%%%%%%%%%%%%%%%%%%%%

The beta function is

\[
\boxed{
B(\alpha,\beta)
=
\int_0^1
x^{\alpha-1}(1-x)^{\beta-1}\,dx.
}
\]

It is related to the gamma function by

\[
\boxed{
B(\alpha,\beta)
=
\frac{
\Gamma(\alpha)\Gamma(\beta)
}{
\Gamma(\alpha+\beta)
}.
}
\]

This identity links beta and gamma distributions algebraically.

%%%%%%%%%%%%%%%%%%%%%%%%%%%%%%%%%%%%%%%%%%%%%%%%%%%%%%%%%%%%
\subsection{Shapes of Beta Distributions}
\label{subsec:beta-shapes}
%%%%%%%%%%%%%%%%%%%%%%%%%%%%%%%%%%%%%%%%%%%%%%%%%%%%%%%%%%%%

The beta family can represent many different shapes.

For example:

\[
\alpha=\beta=1
\]

gives a uniform density on \([0,1]\).

If

\[
\alpha>1,
\qquad
\beta>1,
\]

the density is typically concentrated away from the endpoints.

If one parameter is less than \(1\), the density may rise sharply near an
endpoint.

%%%%%%%%%%%%%%%%%%%%%%%%%%%%%%%%%%%%%%%%%%%%%%%%%%%%%%%%%%%%
\subsection{Uniform Distribution as a Beta Special Case}
\label{subsec:uniform-beta-special-case}
%%%%%%%%%%%%%%%%%%%%%%%%%%%%%%%%%%%%%%%%%%%%%%%%%%%%%%%%%%%%

For

\[
\alpha=\beta=1,
\]

we have

\[
B(1,1)=1.
\]

Therefore,

\[
f_X(x)=1
\]

for

\[
0<x<1.
\]

Hence,

\[
\boxed{
\operatorname{Beta}(1,1)
=
\operatorname{Uniform}(0,1).
}
\]

%%%%%%%%%%%%%%%%%%%%%%%%%%%%%%%%%%%%%%%%%%%%%%%%%%%%%%%%%%%%
\subsection{Weibull Distribution}
\label{subsec:weibull-distribution}
%%%%%%%%%%%%%%%%%%%%%%%%%%%%%%%%%%%%%%%%%%%%%%%%%%%%%%%%%%%%

The Weibull distribution is widely used for lifetime and reliability
models.

\begin{definitionbox}{Weibull Distribution}
A common shape-scale parameterization is

\[
\boxed{
f_X(x)
=
\frac{k}{\lambda}
\left(
\frac{x}{\lambda}
\right)^{k-1}
\exp
\left[
-
\left(
\frac{x}{\lambda}
\right)^k
\right],
\qquad
x\geq0,
}
\]

where

\[
k>0
\]

is the shape parameter and

\[
\lambda>0
\]

is the scale parameter.
\end{definitionbox}

%%%%%%%%%%%%%%%%%%%%%%%%%%%%%%%%%%%%%%%%%%%%%%%%%%%%%%%%%%%%
\subsection{Weibull CDF}
\label{subsec:weibull-cdf}
%%%%%%%%%%%%%%%%%%%%%%%%%%%%%%%%%%%%%%%%%%%%%%%%%%%%%%%%%%%%

For

\[
x\geq0,
\]

the Weibull CDF is

\[
\boxed{
F_X(x)
=
1-
\exp
\left[
-
\left(
\frac{x}{\lambda}
\right)^k
\right].
}
\]

Therefore,

\[
\boxed{
S_X(x)
=
\exp
\left[
-
\left(
\frac{x}{\lambda}
\right)^k
\right].
}
\]

%%%%%%%%%%%%%%%%%%%%%%%%%%%%%%%%%%%%%%%%%%%%%%%%%%%%%%%%%%%%
\subsection{Weibull Hazard Behavior}
\label{subsec:weibull-hazard}
%%%%%%%%%%%%%%%%%%%%%%%%%%%%%%%%%%%%%%%%%%%%%%%%%%%%%%%%%%%%

The Weibull hazard function is

\[
\boxed{
h(x)
=
\frac{k}{\lambda}
\left(
\frac{x}{\lambda}
\right)^{k-1}.
}
\]

Therefore:

\[
k<1
\]

gives decreasing hazard,

\[
k=1
\]

gives constant hazard,

and

\[
k>1
\]

gives increasing hazard.

This makes the Weibull family highly flexible for reliability analysis.

%%%%%%%%%%%%%%%%%%%%%%%%%%%%%%%%%%%%%%%%%%%%%%%%%%%%%%%%%%%%
\subsection{Exponential as a Weibull Special Case}
\label{subsec:exponential-weibull-special-case}
%%%%%%%%%%%%%%%%%%%%%%%%%%%%%%%%%%%%%%%%%%%%%%%%%%%%%%%%%%%%

Set

\[
k=1.
\]

Then

\[
f_X(x)
=
\frac1\lambda
e^{-x/\lambda}.
\]

This is an exponential distribution with rate

\[
\frac1\lambda.
\]

Therefore,

\[
\boxed{
\text{Weibull shape }1
=
\text{exponential}.
}
\]

%%%%%%%%%%%%%%%%%%%%%%%%%%%%%%%%%%%%%%%%%%%%%%%%%%%%%%%%%%%%
\subsection{Cauchy Distribution}
\label{subsec:cauchy-distribution}
%%%%%%%%%%%%%%%%%%%%%%%%%%%%%%%%%%%%%%%%%%%%%%%%%%%%%%%%%%%%

The Cauchy distribution is an important heavy-tailed distribution.

\begin{definitionbox}{Cauchy Distribution}
A standard Cauchy random variable has density

\[
\boxed{
f_X(x)
=
\frac{1}{\pi(1+x^2)},
\qquad
-\infty<x<\infty.
}
\]
\end{definitionbox}

A location-scale Cauchy distribution has density

\[
\boxed{
f_X(x)
=
\frac{1}{
\pi\gamma
\left[
1+
\left(
\frac{x-x_0}{\gamma}
\right)^2
\right]
},
}
\]

where \(x_0\) is the location and \(\gamma>0\) is the scale.

%%%%%%%%%%%%%%%%%%%%%%%%%%%%%%%%%%%%%%%%%%%%%%%%%%%%%%%%%%%%
\subsection{Cauchy CDF}
\label{subsec:cauchy-cdf}
%%%%%%%%%%%%%%%%%%%%%%%%%%%%%%%%%%%%%%%%%%%%%%%%%%%%%%%%%%%%

For the standard Cauchy distribution,

\[
\boxed{
F_X(x)
=
\frac12
+
\frac1\pi
\arctan(x).
}
\]

Its tails decay much more slowly than normal tails.

This produces several unusual statistical properties that will be
discussed when expectation and moments are introduced.

%%%%%%%%%%%%%%%%%%%%%%%%%%%%%%%%%%%%%%%%%%%%%%%%%%%%%%%%%%%%
\subsection{Heavy Tails}
\label{subsec:heavy-tails}
%%%%%%%%%%%%%%%%%%%%%%%%%%%%%%%%%%%%%%%%%%%%%%%%%%%%%%%%%%%%

A distribution is informally called heavy-tailed if unusually large
values retain more probability than they would under rapidly decaying
models such as the normal distribution.

Heavy tails matter because extreme observations may occur much more
frequently than normal modeling suggests.

Examples include:

\[
\boxed{
\text{Cauchy}
}
\]

and

\[
\boxed{
\text{Pareto}.
}
\]

%%%%%%%%%%%%%%%%%%%%%%%%%%%%%%%%%%%%%%%%%%%%%%%%%%%%%%%%%%%%
\subsection{Pareto Distribution}
\label{subsec:pareto-distribution}
%%%%%%%%%%%%%%%%%%%%%%%%%%%%%%%%%%%%%%%%%%%%%%%%%%%%%%%%%%%%

The Pareto distribution is a classical power-law model.

\begin{definitionbox}{Pareto Distribution}
For scale parameter \(x_m>0\) and shape parameter \(\alpha>0\),

\[
\boxed{
f_X(x)
=
\frac{
\alpha x_m^\alpha
}{
x^{\alpha+1}
},
\qquad
x\geq x_m.
}
\]
\end{definitionbox}

Its CDF is

\[
\boxed{
F_X(x)
=
1-
\left(
\frac{x_m}{x}
\right)^\alpha,
\qquad
x\geq x_m.
}
\]

%%%%%%%%%%%%%%%%%%%%%%%%%%%%%%%%%%%%%%%%%%%%%%%%%%%%%%%%%%%%
\subsection{Pareto Survival Function}
\label{subsec:pareto-survival}
%%%%%%%%%%%%%%%%%%%%%%%%%%%%%%%%%%%%%%%%%%%%%%%%%%%%%%%%%%%%

For \(x\geq x_m\),

\[
\boxed{
P(X>x)
=
\left(
\frac{x_m}{x}
\right)^\alpha.
}
\]

Unlike exponential tails, Pareto tails decrease according to a power
law.

This makes large values comparatively more likely.

%%%%%%%%%%%%%%%%%%%%%%%%%%%%%%%%%%%%%%%%%%%%%%%%%%%%%%%%%%%%
\subsection{Triangular Distribution}
\label{subsec:triangular-distribution}
%%%%%%%%%%%%%%%%%%%%%%%%%%%%%%%%%%%%%%%%%%%%%%%%%%%%%%%%%%%%

A triangular distribution is useful when only a minimum, maximum, and
most likely value are available.

Let

\[
a<c<b.
\]

A triangular density is

\[
\boxed{
f_X(x)
=
\begin{cases}
\dfrac{2(x-a)}{(b-a)(c-a)},
&
a\leq x\leq c,
\\[3mm]
\dfrac{2(b-x)}{(b-a)(b-c)},
&
c<x\leq b,
\\[3mm]
0,
&
\text{otherwise}.
\end{cases}
}
\]

The value \(c\) is the mode.

%%%%%%%%%%%%%%%%%%%%%%%%%%%%%%%%%%%%%%%%%%%%%%%%%%%%%%%%%%%%
\subsection{Triangular Distribution Applications}
\label{subsec:triangular-applications}
%%%%%%%%%%%%%%%%%%%%%%%%%%%%%%%%%%%%%%%%%%%%%%%%%%%%%%%%%%%%

Triangular models are often used when detailed data are unavailable but
reasonable lower, upper, and most likely estimates exist.

Applications include:

\begin{itemize}
    \item project duration estimation;
    \item engineering uncertainty;
    \item simulation;
    \item preliminary risk analysis.
\end{itemize}

%%%%%%%%%%%%%%%%%%%%%%%%%%%%%%%%%%%%%%%%%%%%%%%%%%%%%%%%%%%%
\subsection{Lognormal Distribution Preview}
\label{subsec:lognormal-distribution-preview}
%%%%%%%%%%%%%%%%%%%%%%%%%%%%%%%%%%%%%%%%%%%%%%%%%%%%%%%%%%%%

Suppose

\[
Y\sim N(\mu,\sigma^2)
\]

and define

\[
X=e^Y.
\]

Then \(X\) is lognormally distributed.

Equivalently,

\[
\boxed{
\ln X
\sim
N(\mu,\sigma^2).
}
\]

The support is

\[
x>0.
\]

Lognormal models are useful for positive variables generated
multiplicatively.

%%%%%%%%%%%%%%%%%%%%%%%%%%%%%%%%%%%%%%%%%%%%%%%%%%%%%%%%%%%%
\subsection{Lognormal Density}
\label{subsec:lognormal-density}
%%%%%%%%%%%%%%%%%%%%%%%%%%%%%%%%%%%%%%%%%%%%%%%%%%%%%%%%%%%%

The lognormal density is

\[
\boxed{
f_X(x)
=
\frac{
1
}{
x\sigma\sqrt{2\pi}
}
\exp
\left[
-
\frac{
(\ln x-\mu)^2
}{
2\sigma^2
}
\right],
\qquad
x>0.
}
\]

The distribution is typically right-skewed.

%%%%%%%%%%%%%%%%%%%%%%%%%%%%%%%%%%%%%%%%%%%%%%%%%%%%%%%%%%%%
\subsection{Extreme-Value Distributions Preview}
\label{subsec:extreme-value-distributions-preview}
%%%%%%%%%%%%%%%%%%%%%%%%%%%%%%%%%%%%%%%%%%%%%%%%%%%%%%%%%%%%

When studying maxima or minima of large samples, special limiting
families called extreme-value distributions arise.

Examples include:

\[
\text{Gumbel},
\qquad
\text{Fr\'echet},
\qquad
\text{Weibull-type extreme-value laws}.
\]

These distributions are important in:

\begin{itemize}
    \item flood modeling;
    \item material failure;
    \item financial risk;
    \item environmental extremes.
\end{itemize}

%%%%%%%%%%%%%%%%%%%%%%%%%%%%%%%%%%%%%%%%%%%%%%%%%%%%%%%%%%%%
\subsection{Continuous Quantiles}
\label{subsec:continuous-quantiles}
%%%%%%%%%%%%%%%%%%%%%%%%%%%%%%%%%%%%%%%%%%%%%%%%%%%%%%%%%%%%

For a continuous distribution with strictly increasing CDF, the
\(p\)-quantile \(q_p\) satisfies

\[
\boxed{
F_X(q_p)=p.
}
\]

Therefore,

\[
\boxed{
q_p=F_X^{-1}(p).
}
\]

This provides a systematic method for computing percentiles.

%%%%%%%%%%%%%%%%%%%%%%%%%%%%%%%%%%%%%%%%%%%%%%%%%%%%%%%%%%%%
\subsection{Worked Example: Exponential Quantile}
\label{subsec:worked-exponential-quantile}
%%%%%%%%%%%%%%%%%%%%%%%%%%%%%%%%%%%%%%%%%%%%%%%%%%%%%%%%%%%%

\begin{workedexamplebox}{Finding an Exponential Percentile}

Let

\[
X\sim\operatorname{Exp}(\lambda).
\]

We want \(q_p\) satisfying

\[
P(X\leq q_p)=p.
\]

Using

\[
F_X(x)=1-e^{-\lambda x},
\]

we solve

\[
1-e^{-\lambda q_p}=p.
\]

Then

\[
e^{-\lambda q_p}=1-p.
\]

Taking logarithms,

\[
-\lambda q_p
=
\ln(1-p).
\]

\begin{answerbox}
\[
\boxed{
q_p
=
-\frac{1}{\lambda}
\ln(1-p).
}
\]
\end{answerbox}

\end{workedexamplebox}

%%%%%%%%%%%%%%%%%%%%%%%%%%%%%%%%%%%%%%%%%%%%%%%%%%%%%%%%%%%%
\subsection{Median of an Exponential Variable}
\label{subsec:exponential-median}
%%%%%%%%%%%%%%%%%%%%%%%%%%%%%%%%%%%%%%%%%%%%%%%%%%%%%%%%%%%%

Set

\[
p=\frac12.
\]

Then

\[
q_{0.5}
=
-\frac1\lambda
\ln\left(\frac12\right).
\]

Since

\[
-\ln\left(\frac12\right)
=
\ln2,
\]

the median is

\[
\boxed{
\frac{\ln2}{\lambda}.
}
\]

%%%%%%%%%%%%%%%%%%%%%%%%%%%%%%%%%%%%%%%%%%%%%%%%%%%%%%%%%%%%
\subsection{Transformation Method Preview}
\label{subsec:continuous-transformation-method-preview}
%%%%%%%%%%%%%%%%%%%%%%%%%%%%%%%%%%%%%%%%%%%%%%%%%%%%%%%%%%%%

Suppose

\[
Y=g(X).
\]

There are two common methods for finding the distribution of \(Y\):

\begin{enumerate}
    \item transform the CDF;
    \item use the density change-of-variables formula.
\end{enumerate}

For a differentiable one-to-one transformation,

\[
\boxed{
f_Y(y)
=
f_X(g^{-1}(y))
\left|
\frac{d}{dy}
g^{-1}(y)
\right|.
}
\]

This is developed fully in Section~\ref{sec:functions-random-variables}.

%%%%%%%%%%%%%%%%%%%%%%%%%%%%%%%%%%%%%%%%%%%%%%%%%%%%%%%%%%%%
\subsection{Worked Example: Scaling a Uniform Variable}
\label{subsec:worked-uniform-scaling}
%%%%%%%%%%%%%%%%%%%%%%%%%%%%%%%%%%%%%%%%%%%%%%%%%%%%%%%%%%%%

\begin{workedexamplebox}{Transforming a Uniform Random Variable}

Suppose

\[
X\sim\operatorname{Uniform}(0,1)
\]

and

\[
Y=4X+2.
\]

Then

\[
0\leq X\leq1
\]

maps to

\[
2\leq Y\leq6.
\]

The transformation stretches the interval length from \(1\) to \(4\).

Therefore the density height becomes

\[
\frac14.
\]

\begin{answerbox}
\[
\boxed{
Y\sim\operatorname{Uniform}(2,6).
}
\]
\end{answerbox}

\end{workedexamplebox}

%%%%%%%%%%%%%%%%%%%%%%%%%%%%%%%%%%%%%%%%%%%%%%%%%%%%%%%%%%%%
\subsection{Location-Scale Families}
\label{subsec:location-scale-families}
%%%%%%%%%%%%%%%%%%%%%%%%%%%%%%%%%%%%%%%%%%%%%%%%%%%%%%%%%%%%

Suppose \(Z\) has some fixed distribution and define

\[
\boxed{
X=\mu+\sigma Z,
\qquad
\sigma>0.
}
\]

Then \(X\) belongs to the corresponding location-scale family.

The density transforms as

\[
\boxed{
f_X(x)
=
\frac1{\sigma}
f_Z
\left(
\frac{x-\mu}{\sigma}
\right).
}
\]

Normal and Cauchy distributions are important examples.

%%%%%%%%%%%%%%%%%%%%%%%%%%%%%%%%%%%%%%%%%%%%%%%%%%%%%%%%%%%%
\subsection{Symmetric Continuous Distributions}
\label{subsec:symmetric-continuous-distributions}
%%%%%%%%%%%%%%%%%%%%%%%%%%%%%%%%%%%%%%%%%%%%%%%%%%%%%%%%%%%%

A density is symmetric about \(c\) if

\[
\boxed{
f(c-x)=f(c+x)
}
\]

whenever both sides are defined.

For a symmetric distribution,

\[
c
\]

acts as a natural center.

Examples include:

\[
\boxed{
\text{normal},
\qquad
\text{Cauchy},
\qquad
\text{symmetric uniform}.
}
\]

%%%%%%%%%%%%%%%%%%%%%%%%%%%%%%%%%%%%%%%%%%%%%%%%%%%%%%%%%%%%
\subsection{Skewed Distributions}
\label{subsec:skewed-continuous-distributions}
%%%%%%%%%%%%%%%%%%%%%%%%%%%%%%%%%%%%%%%%%%%%%%%%%%%%%%%%%%%%

A distribution is right-skewed when it has a long right tail.

Examples include many:

\[
\text{gamma},
\qquad
\text{Weibull},
\qquad
\text{lognormal},
\qquad
\text{Pareto}
\]

models.

A left-skewed distribution has a longer left tail.

Skewness will be quantified using moments later.

%%%%%%%%%%%%%%%%%%%%%%%%%%%%%%%%%%%%%%%%%%%%%%%%%%%%%%%%%%%%
\subsection{Support and Density Shape}
\label{subsec:support-density-shape}
%%%%%%%%%%%%%%%%%%%%%%%%%%%%%%%%%%%%%%%%%%%%%%%%%%%%%%%%%%%%

The support provides an important clue when choosing a distribution.

For example:

\[
\operatorname{Normal}
\]

has support

\[
(-\infty,\infty).
\]

\[
\operatorname{Exponential},
\operatorname{Gamma},
\operatorname{Weibull}
\]

have support on

\[
[0,\infty).
\]

\[
\operatorname{Beta}
\]

has support on

\[
[0,1].
\]

\[
\operatorname{Pareto}
\]

has support on

\[
[x_m,\infty).
\]

%%%%%%%%%%%%%%%%%%%%%%%%%%%%%%%%%%%%%%%%%%%%%%%%%%%%%%%%%%%%
\subsection{Continuous Distribution Selection Table}
\label{subsec:continuous-distribution-selection-table}
%%%%%%%%%%%%%%%%%%%%%%%%%%%%%%%%%%%%%%%%%%%%%%%%%%%%%%%%%%%%

\[
\begin{array}{c|c}
\text{Problem Structure}
&
\text{Possible Distribution}
\\
\hline
\text{Equal density on a finite interval}
&
\operatorname{Uniform}
\\[1mm]
\text{Memoryless waiting time}
&
\operatorname{Exponential}
\\[1mm]
\text{Symmetric bell-shaped measurement}
&
\operatorname{Normal}
\\[1mm]
\text{Positive waiting time to several events}
&
\operatorname{Gamma}
\\[1mm]
\text{Random proportion}
&
\operatorname{Beta}
\\[1mm]
\text{Flexible lifetime model}
&
\operatorname{Weibull}
\\[1mm]
\text{Positive multiplicative quantity}
&
\operatorname{Lognormal}
\\[1mm]
\text{Heavy-tailed positive quantity}
&
\operatorname{Pareto}
\end{array}
\]

%%%%%%%%%%%%%%%%%%%%%%%%%%%%%%%%%%%%%%%%%%%%%%%%%%%%%%%%%%%%
\subsection{Continuous Distribution Relationships}
\label{subsec:continuous-distribution-relationships}
%%%%%%%%%%%%%%%%%%%%%%%%%%%%%%%%%%%%%%%%%%%%%%%%%%%%%%%%%%%%

Several important families are nested.

\[
\boxed{
\operatorname{Exp}(\lambda)
=
\operatorname{Gamma}(1,\lambda).
}
\]

Also,

\[
\boxed{
\chi^2_\nu
=
\operatorname{Gamma}
\left(
\frac{\nu}{2},
\frac12
\right).
}
\]

Further,

\[
\boxed{
\operatorname{Beta}(1,1)
=
\operatorname{Uniform}(0,1).
}
\]

And Weibull shape \(1\) reduces to an exponential model.

These relationships help organize the distribution families.

%%%%%%%%%%%%%%%%%%%%%%%%%%%%%%%%%%%%%%%%%%%%%%%%%%%%%%%%%%%%
\subsection{Continuous Versus Discrete Probability}
\label{subsec:continuous-versus-discrete}
%%%%%%%%%%%%%%%%%%%%%%%%%%%%%%%%%%%%%%%%%%%%%%%%%%%%%%%%%%%%

For a discrete random variable:

\[
\boxed{
P(X=x)
}
\]

can be positive.

For an absolutely continuous random variable:

\[
\boxed{
P(X=x)=0.
}
\]

Discrete probabilities are summed:

\[
\boxed{
P(X\in A)
=
\sum_{x\in A}p_X(x).
}
\]

Continuous probabilities are integrated:

\[
\boxed{
P(X\in A)
=
\int_A f_X(x)\,dx.
}
\]

%%%%%%%%%%%%%%%%%%%%%%%%%%%%%%%%%%%%%%%%%%%%%%%%%%%%%%%%%%%%
\subsection{PMF Versus PDF}
\label{subsec:pmf-versus-pdf}
%%%%%%%%%%%%%%%%%%%%%%%%%%%%%%%%%%%%%%%%%%%%%%%%%%%%%%%%%%%%

\[
\begin{array}{c|c|c}
&
\text{Discrete}
&
\text{Continuous}
\\
\hline
\text{Basic Function}
&
p_X(x)
&
f_X(x)
\\[1mm]
\text{Point Probability}
&
P(X=x)=p_X(x)
&
P(X=x)=0
\\[1mm]
\text{Interval Probability}
&
\displaystyle\sum p_X(x)
&
\displaystyle\int f_X(x)\,dx
\\[3mm]
\text{Normalization}
&
\displaystyle\sum_xp_X(x)=1
&
\displaystyle\int_{-\infty}^{\infty}f_X(x)\,dx=1
\end{array}
\]

%%%%%%%%%%%%%%%%%%%%%%%%%%%%%%%%%%%%%%%%%%%%%%%%%%%%%%%%%%%%
\subsection{Improper Integrals in Probability}
\label{subsec:improper-integrals-probability}
%%%%%%%%%%%%%%%%%%%%%%%%%%%%%%%%%%%%%%%%%%%%%%%%%%%%%%%%%%%%

Many continuous densities have infinite support.

For example,

\[
\operatorname{Normal}
\]

has support on the entire real line.

Therefore normalization involves an improper integral:

\[
\boxed{
\int_{-\infty}^{\infty}
f_X(x)\,dx.
}
\]

Similarly, positive lifetime distributions often require

\[
\boxed{
\int_0^\infty
f_X(x)\,dx.
}
\]

Convergence of these integrals is essential.

%%%%%%%%%%%%%%%%%%%%%%%%%%%%%%%%%%%%%%%%%%%%%%%%%%%%%%%%%%%%
\subsection{The Gaussian Integral}
\label{subsec:gaussian-integral}
%%%%%%%%%%%%%%%%%%%%%%%%%%%%%%%%%%%%%%%%%%%%%%%%%%%%%%%%%%%%

Normalization of the standard normal density relies on the famous
Gaussian integral

\[
\boxed{
\int_{-\infty}^{\infty}
e^{-x^2/2}\,dx
=
\sqrt{2\pi}.
}
\]

Therefore,

\[
\boxed{
\int_{-\infty}^{\infty}
\frac1{\sqrt{2\pi}}
e^{-x^2/2}\,dx
=
1.
}
\]

A standard proof squares the integral and converts to polar coordinates.

%%%%%%%%%%%%%%%%%%%%%%%%%%%%%%%%%%%%%%%%%%%%%%%%%%%%%%%%%%%%
\subsection{Worked Example: Gaussian Integral Outline}
\label{subsec:worked-gaussian-integral-outline}
%%%%%%%%%%%%%%%%%%%%%%%%%%%%%%%%%%%%%%%%%%%%%%%%%%%%%%%%%%%%

\begin{workedexamplebox}{Why the Normalizing Constant Works}

Let

\[
I
=
\int_{-\infty}^{\infty}
e^{-x^2/2}\,dx.
\]

Then

\[
I^2
=
\int_{\mathbb{R}^2}
e^{-(x^2+y^2)/2}\,dx\,dy.
\]

Using polar coordinates,

\[
I^2
=
\int_0^{2\pi}
\int_0^\infty
e^{-r^2/2}
r\,dr\,d\theta.
\]

The radial integral equals \(1\), so

\[
I^2=2\pi.
\]

Hence,

\begin{answerbox}
\[
\boxed{
I=\sqrt{2\pi}.
}
\]
\end{answerbox}

\end{workedexamplebox}

%%%%%%%%%%%%%%%%%%%%%%%%%%%%%%%%%%%%%%%%%%%%%%%%%%%%%%%%%%%%
\subsection{Probability Integral Transform Preview}
\label{subsec:probability-integral-transform-preview}
%%%%%%%%%%%%%%%%%%%%%%%%%%%%%%%%%%%%%%%%%%%%%%%%%%%%%%%%%%%%

Suppose \(X\) has a continuous strictly increasing CDF \(F_X\).

Then the transformed variable

\[
\boxed{
U=F_X(X)
}
\]

is uniformly distributed on

\[
(0,1).
\]

Conversely, if

\[
U\sim\operatorname{Uniform}(0,1),
\]

then

\[
\boxed{
X=F_X^{-1}(U)
}
\]

has CDF \(F_X\).

This provides a foundation for random-number simulation.

%%%%%%%%%%%%%%%%%%%%%%%%%%%%%%%%%%%%%%%%%%%%%%%%%%%%%%%%%%%%
\subsection{Inverse Transform Sampling}
\label{subsec:inverse-transform-sampling}
%%%%%%%%%%%%%%%%%%%%%%%%%%%%%%%%%%%%%%%%%%%%%%%%%%%%%%%%%%%%

To simulate a continuous random variable with CDF \(F\):

\begin{enumerate}
    \item generate
    \[
    U\sim\operatorname{Uniform}(0,1);
    \]

    \item set
    \[
    X=F^{-1}(U).
    \]
\end{enumerate}

Then \(X\) has the desired distribution under appropriate regularity.

%%%%%%%%%%%%%%%%%%%%%%%%%%%%%%%%%%%%%%%%%%%%%%%%%%%%%%%%%%%%
\subsection{Worked Example: Simulating an Exponential Variable}
\label{subsec:worked-inverse-exponential}
%%%%%%%%%%%%%%%%%%%%%%%%%%%%%%%%%%%%%%%%%%%%%%%%%%%%%%%%%%%%

\begin{workedexamplebox}{Inverse Transform for the Exponential Law}

For

\[
X\sim\operatorname{Exp}(\lambda),
\]

the CDF is

\[
F(x)=1-e^{-\lambda x}.
\]

Set

\[
U=F(X).
\]

Solve:

\[
U
=
1-e^{-\lambda X}.
\]

Then

\[
e^{-\lambda X}=1-U.
\]

Thus,

\[
X
=
-\frac1\lambda
\ln(1-U).
\]

Since \(1-U\) is also uniform on \((0,1)\), one often writes

\begin{answerbox}
\[
\boxed{
X
=
-\frac1\lambda
\ln U.
}
\]
\end{answerbox}

\end{workedexamplebox}

%%%%%%%%%%%%%%%%%%%%%%%%%%%%%%%%%%%%%%%%%%%%%%%%%%%%%%%%%%%%
\subsection{Continuous Model Selection Workflow}
\label{subsec:continuous-model-selection-workflow}
%%%%%%%%%%%%%%%%%%%%%%%%%%%%%%%%%%%%%%%%%%%%%%%%%%%%%%%%%%%%

When selecting a continuous model, ask:

\begin{enumerate}
    \item What values can the variable take?
    \item Is the support bounded or unbounded?
    \item Must values be positive?
    \item Is the distribution symmetric or skewed?
    \item Is a memoryless assumption reasonable?
    \item Are heavy tails plausible?
    \item Does the mechanism involve waiting time?
    \item Is the variable a proportion?
    \item Is the model intended for measurements, lifetimes, or extreme
          values?
\end{enumerate}

The support and physical mechanism are often more informative than the
visual shape alone.

%%%%%%%%%%%%%%%%%%%%%%%%%%%%%%%%%%%%%%%%%%%%%%%%%%%%%%%%%%%%
\subsection{Common Errors}
\label{subsec:continuous-distributions-common-errors}
%%%%%%%%%%%%%%%%%%%%%%%%%%%%%%%%%%%%%%%%%%%%%%%%%%%%%%%%%%%%

\subsubsection{Treating a Density Value as a Probability}

In general,

\[
\boxed{
f_X(x)\neq P(X=x).
}
\]

For continuous variables,

\[
P(X=x)=0.
\]

%%%%%%%%%%%%%%%%%%%%%%%%%%%%%%%%%%%%%%%%%%%%%%%%%%%%%%%%%%%%
\subsubsection{Assuming a PDF Must Be Less Than One}

A density can exceed \(1\).

The requirement is that the total area equals \(1\).

%%%%%%%%%%%%%%%%%%%%%%%%%%%%%%%%%%%%%%%%%%%%%%%%%%%%%%%%%%%%
\subsubsection{Forgetting to Check Normalization}

A nonnegative function is not automatically a PDF.

It must also satisfy

\[
\int_{-\infty}^{\infty}f(x)\,dx=1.
\]

%%%%%%%%%%%%%%%%%%%%%%%%%%%%%%%%%%%%%%%%%%%%%%%%%%%%%%%%%%%%
\subsubsection{Integrating Outside the Support Incorrectly}

The density is zero outside its support.

Probability integrals should respect the actual support of the random
variable.

%%%%%%%%%%%%%%%%%%%%%%%%%%%%%%%%%%%%%%%%%%%%%%%%%%%%%%%%%%%%
\subsubsection{Confusing Rate and Scale Parameters}

For an exponential rate parameterization,

\[
f(x)=\lambda e^{-\lambda x}.
\]

Some sources instead use a scale parameter

\[
\theta=\frac1\lambda.
\]

Always check the convention.

%%%%%%%%%%%%%%%%%%%%%%%%%%%%%%%%%%%%%%%%%%%%%%%%%%%%%%%%%%%%
\subsubsection{Confusing Gamma Rate and Scale}

A gamma density may be parameterized using rate \(\lambda\) or scale
\(\theta\).

These satisfy

\[
\boxed{
\theta=\frac1\lambda.
}
\]

The formulas must not be mixed.

%%%%%%%%%%%%%%%%%%%%%%%%%%%%%%%%%%%%%%%%%%%%%%%%%%%%%%%%%%%%
\subsubsection{Using Variance Instead of Standard Deviation in a \(z\)-Score}

For

\[
X\sim N(\mu,\sigma^2),
\]

standardization uses

\[
\boxed{
Z=\frac{X-\mu}{\sigma},
}
\]

not division by \(\sigma^2\).

%%%%%%%%%%%%%%%%%%%%%%%%%%%%%%%%%%%%%%%%%%%%%%%%%%%%%%%%%%%%
\subsubsection{Forgetting Normal Tail Complements}

If

\[
P(Z\leq z)=\Phi(z),
\]

then

\[
\boxed{
P(Z>z)=1-\Phi(z).
}
\]

%%%%%%%%%%%%%%%%%%%%%%%%%%%%%%%%%%%%%%%%%%%%%%%%%%%%%%%%%%%%
\subsubsection{Ignoring Normal Symmetry}

For the standard normal distribution,

\[
\boxed{
\Phi(-z)=1-\Phi(z).
}
\]

This often avoids unnecessary computations.

%%%%%%%%%%%%%%%%%%%%%%%%%%%%%%%%%%%%%%%%%%%%%%%%%%%%%%%%%%%%
\subsubsection{Using an Exponential Model When Hazard Is Not Constant}

The exponential model assumes a constant hazard rate.

If failure rate changes with age, a Weibull or other lifetime model may
be more appropriate.

%%%%%%%%%%%%%%%%%%%%%%%%%%%%%%%%%%%%%%%%%%%%%%%%%%%%%%%%%%%%
\subsubsection{Assuming All Positive Variables Are Normally Distributed}

The normal distribution allows negative values.

For strictly positive and strongly skewed quantities, gamma, Weibull,
lognormal, or Pareto models may be more appropriate.

%%%%%%%%%%%%%%%%%%%%%%%%%%%%%%%%%%%%%%%%%%%%%%%%%%%%%%%%%%%%
\subsection{Applications}
\label{subsec:continuous-distributions-applications}
%%%%%%%%%%%%%%%%%%%%%%%%%%%%%%%%%%%%%%%%%%%%%%%%%%%%%%%%%%%%

\begin{applicationbox}

\textbf{Measurement Science.}
Normal distributions often model measurement variation and aggregated
small errors.

\medskip

\textbf{Reliability Engineering.}
Exponential and Weibull distributions model component lifetimes and
failure behavior.

\medskip

\textbf{Queueing Theory.}
Exponential waiting times arise naturally in Poisson-process models.

\medskip

\textbf{Statistics.}
Normal, chi-square, beta, gamma, and related distributions appear
throughout statistical inference.

\medskip

\textbf{Finance.}
Lognormal and heavy-tailed distributions are used in models of prices,
returns, and risk.

\medskip

\textbf{Environmental Science.}
Gamma, Weibull, and extreme-value models are used for rainfall,
wind, floods, and other environmental quantities.

\medskip

\textbf{Biology.}
Continuous distributions model lifetimes, concentrations, biological
measurements, and rates.

\medskip

\textbf{Simulation.}
Uniform random variables are transformed into many other distributions
using inverse-CDF and related methods.

\medskip

\textbf{Project Management.}
Triangular distributions model uncertain durations when only minimum,
maximum, and most likely values are known.

\end{applicationbox}

%%%%%%%%%%%%%%%%%%%%%%%%%%%%%%%%%%%%%%%%%%%%%%%%%%%%%%%%%%%%
\subsection{Exercises}
\label{subsec:continuous-distributions-exercises}
%%%%%%%%%%%%%%%%%%%%%%%%%%%%%%%%%%%%%%%%%%%%%%%%%%%%%%%%%%%%

\begin{exercise}[1]
State the two conditions required for a valid probability density
function.
\end{exercise}

\begin{exercise}[1]
Explain why

\[
P(X=x)=0
\]

for a continuous random variable.
\end{exercise}

\begin{exercise}[1]
Suppose \(X\) has density

\[
f(x)
=
\begin{cases}
c,&0\leq x\leq4,\\
0,&\text{otherwise}.
\end{cases}
\]

Find \(c\).
\end{exercise}

\begin{exercise}[1]
Identify the distribution in the previous exercise.
\end{exercise}

\begin{exercise}[1]
State the PDF of

\[
\operatorname{Uniform}(a,b).
\]
\end{exercise}

\begin{exercise}[1]
State the PDF of

\[
\operatorname{Exp}(\lambda).
\]
\end{exercise}

\begin{exercise}[1]
State the CDF of

\[
\operatorname{Exp}(\lambda).
\]
\end{exercise}

\begin{exercise}[1]
State the PDF of

\[
N(\mu,\sigma^2).
\]
\end{exercise}

\begin{exercise}[1]
Define the standard normal distribution.
\end{exercise}

\begin{exercise}[1]
State the standardization formula for a normal random variable.
\end{exercise}

\begin{exercise}[2]
Suppose

\[
X\sim\operatorname{Uniform}(2,10).
\]

Find

\[
P(4\leq X\leq7).
\]
\end{exercise}

\begin{exercise}[2]
Find the CDF of

\[
X\sim\operatorname{Uniform}(2,10).
\]
\end{exercise}

\begin{exercise}[2]
Suppose

\[
X\sim\operatorname{Exp}(0.5).
\]

Find

\[
P(X>4).
\]
\end{exercise}

\begin{exercise}[2]
For the same variable, find

\[
P(X\leq4).
\]
\end{exercise}

\begin{exercise}[2]
Suppose

\[
X\sim\operatorname{Exp}(3).
\]

Find the median.
\end{exercise}

\begin{exercise}[2]
Suppose

\[
X\sim N(20,4^2).
\]

Convert

\[
x=28
\]

to a standard-normal \(z\)-score.
\end{exercise}

\begin{exercise}[2]
For the same variable, express

\[
P(X\leq28)
\]

in terms of \(\Phi\).
\end{exercise}

\begin{exercise}[2]
Express

\[
P(16<X<24)
\]

for

\[
X\sim N(20,4^2)
\]

in terms of \(\Phi\).
\end{exercise}

\begin{exercise}[2]
Use symmetry to simplify

\[
P(-1.5<Z<1.5)
\]

for \(Z\sim N(0,1)\).
\end{exercise}

\begin{exercise}[2]
Find the \(p\)-quantile of

\[
\operatorname{Exp}(\lambda).
\]
\end{exercise}

\begin{exercise}[2]
Show that

\[
\operatorname{Beta}(1,1)
\]

is uniform on \([0,1]\).
\end{exercise}

\begin{exercise}[2]
Show that

\[
\operatorname{Gamma}(1,\lambda)
\]

is exponential with rate \(\lambda\).
\end{exercise}

\begin{exercise}[2]
Identify the gamma parameters corresponding to

\[
\chi^2_8.
\]
\end{exercise}

\begin{exercise}[2]
Show that the Weibull distribution with shape \(1\) is exponential.
\end{exercise}

\begin{exercise}[3]
Find \(c\) such that

\[
f(x)
=
\begin{cases}
cx^2,&0\leq x\leq2,\\
0,&\text{otherwise}
\end{cases}
\]

is a valid PDF.
\end{exercise}

\begin{exercise}[3]
For the previous density, derive the complete CDF.
\end{exercise}

\begin{exercise}[3]
For the same distribution, find

\[
P\left(
\frac12<X<\frac32
\right).
\]
\end{exercise}

\begin{exercise}[3]
Prove the memoryless property of the exponential distribution.
\end{exercise}

\begin{exercise}[3]
Show that the exponential hazard function is constant.
\end{exercise}

\begin{exercise}[3]
Derive the standard normal transformation

\[
Z=\frac{X-\mu}{\sigma}.
\]
\end{exercise}

\begin{exercise}[3]
Show that

\[
\Phi(-z)=1-\Phi(z).
\]
\end{exercise}

\begin{exercise}[3]
Prove the Gaussian integral

\[
\int_{-\infty}^{\infty}
e^{-x^2/2}\,dx
=
\sqrt{2\pi}
\]

using polar coordinates.
\end{exercise}

\begin{exercise}[3]
Verify the gamma recursion

\[
\Gamma(\alpha+1)
=
\alpha\Gamma(\alpha)
\]

using integration by parts.
\end{exercise}

\begin{exercise}[3]
Use the gamma recursion to prove

\[
\Gamma(n)=(n-1)!
\]

for positive integers \(n\).
\end{exercise}

\begin{exercise}[3]
Derive the beta-gamma identity

\[
B(\alpha,\beta)
=
\frac{
\Gamma(\alpha)\Gamma(\beta)
}{
\Gamma(\alpha+\beta)
}.
\]
\end{exercise}

\begin{exercise}[3]
Derive the Weibull CDF from its density.
\end{exercise}

\begin{exercise}[3]
Derive the Pareto survival function.
\end{exercise}

\begin{exercise}[3]
Let

\[
X\sim\operatorname{Uniform}(0,1)
\]

and

\[
Y=3X-1.
\]

Find the density of \(Y\).
\end{exercise}

\begin{exercise}[3]
Let

\[
X\sim\operatorname{Exp}(\lambda)
\]

and

\[
Y=2X.
\]

Find the distribution of \(Y\).
\end{exercise}

\begin{exercise}[3]
Use inverse-transform sampling to derive a simulation formula for an
exponential random variable.
\end{exercise}

\begin{exercise}[4]
Investigate the lognormal distribution as an exponential transform of a
normal random variable.
\end{exercise}

\begin{exercise}[4]
Study Student's \(t\)-distribution and explain how it is constructed
from normal and chi-square variables.
\end{exercise}

\begin{exercise}[4]
Study the \(F\)-distribution and explain its relationship to ratios of
chi-square variables.
\end{exercise}

\begin{exercise}[4]
Investigate generalized extreme-value distributions.
\end{exercise}

\begin{exercise}[4]
Compare gamma, Weibull, and lognormal lifetime models.
\end{exercise}

\begin{exercise}[4]
Study hazard functions for Weibull and Pareto models.
\end{exercise}

\begin{exercise}[4]
Investigate heavy-tailed distributions and compare normal, Cauchy, and
Pareto tail probabilities.
\end{exercise}

\begin{exercise}[4]
Study rejection sampling as an alternative to inverse-transform
sampling.
\end{exercise}

\begin{exercise}[4]
Investigate singular continuous distributions and explain why not every
continuous CDF has a density.
\end{exercise}

\begin{exercise}[4]
Study the role of continuous distributions in survival analysis and
reliability theory.
\end{exercise}

%%%%%%%%%%%%%%%%%%%%%%%%%%%%%%%%%%%%%%%%%%%%%%%%%%%%%%%%%%%%
\subsection{Conceptual Summary}
\label{subsec:continuous-distributions-summary}
%%%%%%%%%%%%%%%%%%%%%%%%%%%%%%%%%%%%%%%%%%%%%%%%%%%%%%%%%%%%

A continuous probability density satisfies

\[
\boxed{
f_X(x)\geq0
}
\]

and

\[
\boxed{
\int_{-\infty}^{\infty}
f_X(x)\,dx
=
1.
}
\]

Probabilities are areas:

\[
\boxed{
P(a\leq X\leq b)
=
\int_a^b f_X(x)\,dx.
}
\]

For absolutely continuous variables,

\[
\boxed{
P(X=x)=0.
}
\]

The CDF is

\[
\boxed{
F_X(x)
=
\int_{-\infty}^{x}
f_X(t)\,dt.
}
\]

When differentiable,

\[
\boxed{
f_X(x)
=
F_X'(x).
}
\]

The continuous uniform distribution has constant density:

\[
\boxed{
f_X(x)
=
\frac1{b-a}.
}
\]

The exponential density is

\[
\boxed{
f_X(x)
=
\lambda e^{-\lambda x},
\qquad
x\geq0.
}
\]

Its survival function is

\[
\boxed{
P(X>x)=e^{-\lambda x}.
}
\]

The normal density is

\[
\boxed{
f_X(x)
=
\frac1{\sigma\sqrt{2\pi}}
\exp
\left(
-\frac{(x-\mu)^2}{2\sigma^2}
\right).
}
\]

Standardization uses

\[
\boxed{
Z
=
\frac{X-\mu}{\sigma}.
}
\]

The standard normal CDF is

\[
\boxed{
\Phi(z)=P(Z\leq z).
}
\]

The gamma family contains the exponential family:

\[
\boxed{
\operatorname{Exp}(\lambda)
=
\operatorname{Gamma}(1,\lambda).
}
\]

The chi-square family satisfies

\[
\boxed{
\chi^2_\nu
=
\operatorname{Gamma}
\left(
\frac{\nu}{2},
\frac12
\right).
}
\]

The beta distribution models random quantities on

\[
[0,1].
\]

The Weibull family provides flexible lifetime models.

The Pareto and Cauchy distributions illustrate heavy-tailed behavior.

Continuous distributions therefore organize recurring models involving

\[
\boxed{
\text{measurement}
+
\text{waiting time}
+
\text{lifetime}
+
\text{proportions}
+
\text{continuous uncertainty}.
}
\]

%%%%%%%%%%%%%%%%%%%%%%%%%%%%%%%%%%%%%%%%%%%%%%%%%%%%%%%%%%%%
\subsection{Section Connections}
\label{subsec:continuous-distributions-connections}
%%%%%%%%%%%%%%%%%%%%%%%%%%%%%%%%%%%%%%%%%%%%%%%%%%%%%%%%%%%%

\chapterconnection{
Continuous Distributions extends the random-variable framework of
Section~\ref{sec:random-variables} from probability masses to densities
and integrals. It complements the Discrete Distributions of
Section~\ref{sec:discrete-distributions}. The next section, Joint
Distributions, studies probability models involving several random
variables simultaneously. Expectation and Moments later provide means,
variances, skewness, and other summaries for the normal, exponential,
gamma, beta, and related families. Functions of Random Variables develops
continuous change-of-variables methods systematically. Probability
Transforms introduces characteristic and moment-generating functions,
while Central Limit Theorems explain the widespread appearance of the
normal distribution. Stochastic Processes later connect Poisson counts
with exponential and gamma waiting times.
}

%%%%%%%%%%%%%%%%%%%%%%%%%%%%%%%%%%%%%%%%%%%%%%%%%%%%%%%%%%%%
% SECTION SOURCE TRACKING
%%%%%%%%%%%%%%%%%%%%%%%%%%%%%%%%%%%%%%%%%%%%%%%%%%%%%%%%%%%%

%------------------------------------------------------------
% 7.4 Continuous Distributions
%------------------------------------------------------------
%
% Source basis:
%   - Standard continuous probability theory.
%   - Standard continuous distribution families.
%   - Standard reliability and waiting-time probability models.
%
% Used for:
%   - probability density functions;
%   - density normalization;
%   - continuous point probabilities;
%   - interval probabilities;
%   - continuous CDFs;
%   - densities from CDFs;
%   - survival functions;
%   - hazard functions preview;
%   - uniform distribution;
%   - exponential distribution;
%   - exponential memorylessness;
%   - exponential-Poisson connection;
%   - normal distribution;
%   - standard normal distribution;
%   - standardization;
%   - z-scores;
%   - normal symmetry;
%   - normal quantiles;
%   - empirical rule;
%   - normal approximation preview;
%   - continuity correction preview;
%   - gamma distribution;
%   - gamma function;
%   - Erlang distribution;
%   - chi-square distribution;
%   - beta distribution;
%   - beta function;
%   - Weibull distribution;
%   - Cauchy distribution;
%   - Pareto distribution;
%   - triangular distribution;
%   - lognormal preview;
%   - extreme-value preview;
%   - quantiles;
%   - inverse-transform sampling;
%   - location-scale families;
%   - support and model selection;
%   - applications and exercises.
%
% Notes:
%   - Means, variances, and higher moments are intentionally deferred to
%     Section 7.6.
%   - Joint continuous distributions appear in Section 7.5.
%   - Transformations and Jacobians are developed more fully in
%     Section 7.7.
%   - Normal approximation is justified later through Central Limit
%     Theorems.
%   - Formal absolute continuity and measure-theoretic distinctions
%     appear in Section 7.11.
%
%%%%%%%%%%%%%%%%%%%%%%%%%%%%%%%%%%%%%%%%%%%%%%%%%%%%%%%%%%%%